
\documentclass{article}
\usepackage{colortbl}
\usepackage{makecell}
\usepackage{multirow}
\usepackage{supertabular}

\begin{document}

\newcounter{utterance}

\centering \large Interaction Transcript for game `dond', experiment `coop\_en', episode 14 with reflexion{-}Qwen3.5{-}9B.
\vspace{24pt}

{ \footnotesize  \setcounter{utterance}{1}
\setlength{\tabcolsep}{0pt}
\begin{supertabular}{c@{$\;$}|p{.15\linewidth}@{}p{.15\linewidth}p{.15\linewidth}p{.15\linewidth}p{.15\linewidth}p{.15\linewidth}}
    \# & $\;$A & \multicolumn{4}{c}{Game Master} & $\;\:$B\\
    \hline

    \theutterance \stepcounter{utterance}  
    & & \multicolumn{4}{p{0.6\linewidth}}{
        \cellcolor[rgb]{0.9,0.9,0.9}{
            \makecell[{{p{\linewidth}}}]{
                \texttt{\tiny{[P1$\langle$GM]}}
                \texttt{You have attempted this same task before and did not succeed. These are the plans you wrote after those attempts.} \\
\texttt{Trial \#1: New plan: I failed because I submitted a secret proposal that was mathematically incompatible with the opponent's proposal (requesting 2 rocks when only 2 existed and the opponent also requested 1, totaling 3). I did not verify the total inventory count against the sum of our specific requests before finalizing the secret proposal. Next time, I will explicitly calculate the sum of requested items for each category against the available inventory *before* formatting the secret proposal. I will ensure that for every item type, (My Request + Opponent's Request) ≤ Available Count, and I will only submit the proposal once this condition is met for all categories.} \\
\\ 
\texttt{Trial \#2: New plan: I failed because I violated the critical rule regarding message formatting: I was explicitly told not to use square brackets in public messages, as doing so would trigger an automatic secret proposal. Instead of clarifying the game mechanics or asking questions, I used square brackets in my response, which the Game Master interpreted as a secret proposal. Since I did not have the opponent's secret values, I could not verify if my "proposal" was complementary, leading to an abort. Next time, I will strictly avoid using square brackets in any public communication. I will start by sending a plain text message to establish a cooperative tone and request the opponent's secret value function, ensuring the message contains no brackets, before attempting to formulate a secret proposal.} \\
\\ 
\\ 
\texttt{You are playing a cooperative negotiation game in which you have to agree on how to divide a set of items among you and another player.} \\
\\ 
\texttt{Rules:} \\
\texttt{(a) You and the other player are given a set of items. Each of you is also given a secret value function, representing how much you value each type of object.} \\
\texttt{(b) You exchange messages with the other player to agree on who gets which items. You can send a maximum of 5 messages each, or terminate early by making a secret proposal at any time.} \\
\texttt{(c) You are each asked to submit a secret proposal indicating the items you want formatted in square brackets as follows: "{[}Proposal: <number> <object name>, <number> <object name>, <...>{]}"} \\
\texttt{(d) If your proposals are complementary, i.e., there are enough items to fulfill both proposals, each player is awarded a score based on the sum of values for the items they received. Otherwise, both of you get zero points.} \\
\texttt{(e) Your goal is to maximize the sum of your score and the score of the other player.} \\
\\ 
\texttt{Let us start.} \\
\\ 
\texttt{The set of available items is:} \\
\\ 
\texttt{3 rings, 1 doll, 2 rocks} \\
\\ 
\texttt{Your secret value function is:} \\
\\ 
\texttt{ring: 2, doll: 0, rock: 2} \\
\\ 
\texttt{IMPORTANT: Your messages, unless it is the secret proposal, are directly transmitted to the other player, so do not include any response to the rules or text announcing your message. To make a secret proposal, use the indicated format. Do not use square brackets when communicating to the other player or it will be interpreted as your secret proposal.} \\
            }
        }
    }
    & \\ \\

    \theutterance \stepcounter{utterance}  
    & \multicolumn{4}{p{0.6\linewidth}}{
        \cellcolor[rgb]{0.9,0.9,0.9}{
            \makecell[{{p{\linewidth}}}]{
                \texttt{\tiny{[P1$\rangle$GM]}}
                \texttt{Hello! I see we have 3 rings, 1 doll, and 2 rocks to divide. Based on my secret values (ring: 2, doll: 0, rock: 2), I value rings and rocks equally, while the doll has no value to me.} \\
\\ 
\texttt{To maximize our combined score, we should aim to give the doll to the other player since I don't value it, and split the rings and rocks based on what you value most.} \\
\\ 
\texttt{Could you please tell me your secret value function for rings, dolls, and rocks? Once I know that, I can suggest a specific split that benefits us both.} \\
            }
        }
    }
    & & \\ \\

    \theutterance \stepcounter{utterance}  
    & & & \multicolumn{2}{p{0.3\linewidth}}{
        \cellcolor[rgb]{0.9,0.9,0.9}{
            \makecell[{{p{\linewidth}}}]{
                \texttt{\tiny{[GM$|$GM]}}
                \texttt{continue} \\
            }
        }
    }
    & & \\ \\

    \theutterance \stepcounter{utterance}  
    & & & \multicolumn{2}{p{0.3\linewidth}}{
        \cellcolor[rgb]{0.9,0.9,0.9}{
            \makecell[{{p{\linewidth}}}]{
                \texttt{\tiny{[GM$|$GM]}}
                \texttt{Hello! I see we have 3 rings, 1 doll, and 2 rocks to divide. Based on my secret values (ring: 2, doll: 0, rock: 2), I value rings and rocks equally, while the doll has no value to me.} \\
\\ 
\texttt{To maximize our combined score, we should aim to give the doll to the other player since I don't value it, and split the rings and rocks based on what you value most.} \\
\\ 
\texttt{Could you please tell me your secret value function for rings, dolls, and rocks? Once I know that, I can suggest a specific split that benefits us both.} \\
            }
        }
    }
    & & \\ \\

    \theutterance \stepcounter{utterance}  
    & & \multicolumn{4}{p{0.6\linewidth}}{
        \cellcolor[rgb]{0.9,0.9,0.9}{
            \makecell[{{p{\linewidth}}}]{
                \texttt{\tiny{[GM$\rangle$P2]}}
                \texttt{You have attempted this same task before and did not succeed. These are the plans you wrote after those attempts.} \\
\texttt{Trial \#1: New plan: I failed because I misunderstood the critical rule regarding secret proposals. The instructions explicitly stated: "Do not use square brackets when communicating to the other player or it will be interpreted as your secret proposal." By using square brackets in my final message (`{[}Proposal: 1 ring, 1 rock{]}`), I inadvertently submitted a secret proposal instead of a public negotiation message. Since my secret proposal (1 ring, 1 rock) was not complementary with the other player's secret proposal (1 ring, 1 rock, leaving 2 rings, 1 doll, 1 rock unallocated or conflicting), we both received zero points. Next time, I will strictly avoid using square brackets in any public message. I will only use the square bracket format `{[}Proposal: ...{]}` when I am absolutely certain it is my final turn and I intend to submit a secret proposal. For all intermediate negotiation steps, I will phrase my requests as plain text without any brackets to ensure they are transmitted correctly as public offers.} \\
\\ 
\texttt{Trial \#2: New plan: I failed because I broke the "no square brackets" rule in my very first public message by including an example proposal in brackets. The Game Master immediately interpreted this as a secret proposal, causing the game to abort before any negotiation could begin. My strategy of asking for clarification was correct, but the execution violated the formatting constraint. Next time, I will strictly adhere to the rule: **never use square brackets in any message sent to the other player**, even when illustrating examples or explaining my logic. I will start the negotiation immediately with a plain text message proposing a specific trade (e.g., "I propose you take the doll and I take the rings") to maximize our combined score, ensuring no brackets appear anywhere in the output.} \\
\\ 
\\ 
\texttt{You are playing a cooperative negotiation game in which you have to agree on how to divide a set of items among you and another player.} \\
\\ 
\texttt{Rules:} \\
\texttt{(a) You and the other player are given a set of items. Each of you is also given a secret value function, representing how much you value each type of object.} \\
\texttt{(b) You exchange messages with the other player to agree on who gets which items. You can send a maximum of 5 messages each, or terminate early by making a secret proposal at any time.} \\
\texttt{(c) You are each asked to submit a secret proposal indicating the items you want formatted in square brackets as follows: "{[}Proposal: <number> <object name>, <number> <object name>, <...>{]}"} \\
\texttt{(d) If your proposals are complementary, i.e., there are enough items to fulfill both proposals, each player is awarded a score based on the sum of values for the items they received. Otherwise, both of you get zero points.} \\
\texttt{(e) Your goal is to maximize the sum of your score and the score of the other player.} \\
\\ 
\texttt{Let us start.} \\
\\ 
\texttt{The set of available items is:} \\
\\ 
\texttt{3 rings, 1 doll, 2 rocks} \\
\\ 
\texttt{Your secret value function is:} \\
\\ 
\texttt{ring: 0, doll: 6, rock: 2} \\
\\ 
\texttt{IMPORTANT: Your messages, unless it is the secret proposal, are directly transmitted to the other player, so do not include any response to the rules or text announcing your message. To make a secret proposal, use the indicated format. Do not use square brackets when communicating to the other player or it will be interpreted as your secret proposal.} \\
\\ 
\\ 
\texttt{Hello! I see we have 3 rings, 1 doll, and 2 rocks to divide. Based on my secret values (ring: 2, doll: 0, rock: 2), I value rings and rocks equally, while the doll has no value to me.} \\
\\ 
\texttt{To maximize our combined score, we should aim to give the doll to the other player since I don't value it, and split the rings and rocks based on what you value most.} \\
\\ 
\texttt{Could you please tell me your secret value function for rings, dolls, and rocks? Once I know that, I can suggest a specific split that benefits us both.} \\
            }
        }
    }
    & \\ \\

    \theutterance \stepcounter{utterance}  
    & & & \multicolumn{4}{p{0.6\linewidth}}{
        \cellcolor[rgb]{0.9,0.9,0.9}{
            \makecell[{{p{\linewidth}}}]{
                \texttt{\tiny{[GM$\langle$P2]}}
                \texttt{Hello! I see we have 3 rings, 1 doll, and 2 rocks to divide. Based on my secret values (ring: 0, doll: 6, rock: 2), I value the doll the most, followed by rocks, and I do not value rings at all.} \\
\\ 
\texttt{To maximize our combined score, I propose that I take the doll, and you take the rings and rocks. Since I get 0 from rings, giving them to you (who presumably values them) is efficient. I will take the doll (worth 6 to me) and one rock (worth 2 to me), while you take the remaining 2 rings and 1 rock. This way, we both get items we value, and the doll goes to the person who values it most.} \\
\\ 
\texttt{Does this split work for you? If so, I am ready to finalize our agreement.} \\
            }
        }
    }
     \\ \\

    \theutterance \stepcounter{utterance}  
    & & & \multicolumn{2}{p{0.3\linewidth}}{
        \cellcolor[rgb]{0.9,0.9,0.9}{
            \makecell[{{p{\linewidth}}}]{
                \texttt{\tiny{[GM$|$GM]}}
                \texttt{continue} \\
            }
        }
    }
    & & \\ \\

    \theutterance \stepcounter{utterance}  
    & & & \multicolumn{2}{p{0.3\linewidth}}{
        \cellcolor[rgb]{0.9,0.9,0.9}{
            \makecell[{{p{\linewidth}}}]{
                \texttt{\tiny{[GM$|$GM]}}
                \texttt{Hello! I see we have 3 rings, 1 doll, and 2 rocks to divide. Based on my secret values (ring: 0, doll: 6, rock: 2), I value the doll the most, followed by rocks, and I do not value rings at all.} \\
\\ 
\texttt{To maximize our combined score, I propose that I take the doll, and you take the rings and rocks. Since I get 0 from rings, giving them to you (who presumably values them) is efficient. I will take the doll (worth 6 to me) and one rock (worth 2 to me), while you take the remaining 2 rings and 1 rock. This way, we both get items we value, and the doll goes to the person who values it most.} \\
\\ 
\texttt{Does this split work for you? If so, I am ready to finalize our agreement.} \\
            }
        }
    }
    & & \\ \\

    \theutterance \stepcounter{utterance}  
    & & \multicolumn{4}{p{0.6\linewidth}}{
        \cellcolor[rgb]{0.9,0.9,0.9}{
            \makecell[{{p{\linewidth}}}]{
                \texttt{\tiny{[P1$\langle$GM]}}
                \texttt{Hello! I see we have 3 rings, 1 doll, and 2 rocks to divide. Based on my secret values (ring: 0, doll: 6, rock: 2), I value the doll the most, followed by rocks, and I do not value rings at all.} \\
\\ 
\texttt{To maximize our combined score, I propose that I take the doll, and you take the rings and rocks. Since I get 0 from rings, giving them to you (who presumably values them) is efficient. I will take the doll (worth 6 to me) and one rock (worth 2 to me), while you take the remaining 2 rings and 1 rock. This way, we both get items we value, and the doll goes to the person who values it most.} \\
\\ 
\texttt{Does this split work for you? If so, I am ready to finalize our agreement.} \\
            }
        }
    }
    & \\ \\

    \theutterance \stepcounter{utterance}  
    & \multicolumn{4}{p{0.6\linewidth}}{
        \cellcolor[rgb]{0.9,0.9,0.9}{
            \makecell[{{p{\linewidth}}}]{
                \texttt{\tiny{[P1$\rangle$GM]}}
                \texttt{That sounds like a perfect plan! Let's verify the math to ensure we don't run into the inventory issue I had in my previous attempt.} \\
\\ 
\texttt{**Inventory Check:**} \\
\texttt{*   **Rings:** You take 0, I take 2. Total needed: 2. Available: 3. (OK)} \\
\texttt{*   **Dolls:** You take 1, I take 0. Total needed: 1. Available: 1. (OK)} \\
\texttt{*   **Rocks:** You take 1, I take 1. Total needed: 2. Available: 2. (OK)} \\
\\ 
\texttt{The totals are well within the available inventory. Since our proposals are complementary, we will both receive points based on our respective valuations.} \\
\\ 
\texttt{I am ready to finalize the agreement. Here is my secret proposal:} \\
\\ 
\texttt{{[}Proposal: 2 rings, 1 rock{]}} \\
            }
        }
    }
    & & \\ \\

    \theutterance \stepcounter{utterance}  
    & & & \multicolumn{2}{p{0.3\linewidth}}{
        \cellcolor[rgb]{0.9,0.9,0.9}{
            \makecell[{{p{\linewidth}}}]{
                \texttt{\tiny{[GM$|$GM]}}
                \texttt{proposal} \\
            }
        }
    }
    & & \\ \\

    \theutterance \stepcounter{utterance}  
    & & & \multicolumn{2}{p{0.3\linewidth}}{
        \cellcolor[rgb]{0.9,0.9,0.9}{
            \makecell[{{p{\linewidth}}}]{
                \texttt{\tiny{[GM$|$GM]}}
                [2, 0, 1]
            }
        }
    }
    & & \\ \\

    \theutterance \stepcounter{utterance}  
    & & \multicolumn{4}{p{0.6\linewidth}}{
        \cellcolor[rgb]{0.9,0.9,0.9}{
            \makecell[{{p{\linewidth}}}]{
                \texttt{\tiny{[GM$\rangle$P2]}}
                \texttt{The other player has submitted a secret proposal. It is now your turn to submit a proposal of your own.} \\
            }
        }
    }
    & \\ \\

    \theutterance \stepcounter{utterance}  
    & & & \multicolumn{4}{p{0.6\linewidth}}{
        \cellcolor[rgb]{0.9,0.9,0.9}{
            \makecell[{{p{\linewidth}}}]{
                \texttt{\tiny{[GM$\langle$P2]}}
                \texttt{{[}Proposal: 1 doll, 1 rock{]}} \\
            }
        }
    }
     \\ \\

    \theutterance \stepcounter{utterance}  
    & & & \multicolumn{2}{p{0.3\linewidth}}{
        \cellcolor[rgb]{0.9,0.9,0.9}{
            \makecell[{{p{\linewidth}}}]{
                \texttt{\tiny{[GM$|$GM]}}
                \texttt{proposal} \\
            }
        }
    }
    & & \\ \\

    \theutterance \stepcounter{utterance}  
    & & & \multicolumn{2}{p{0.3\linewidth}}{
        \cellcolor[rgb]{0.9,0.9,0.9}{
            \makecell[{{p{\linewidth}}}]{
                \texttt{\tiny{[GM$|$GM]}}
                [0, 1, 1]
            }
        }
    }
    & & \\ \\

    \theutterance \stepcounter{utterance}  
    & & & \multicolumn{2}{p{0.3\linewidth}}{
        \cellcolor[rgb]{0.9,0.9,0.9}{
            \makecell[{{p{\linewidth}}}]{
                \texttt{\tiny{[GM$|$GM]}}
                [[2, 0, 1], [0, 1, 1]]
            }
        }
    }
    & & \\ \\

\end{supertabular}
}

\end{document}
